%% supplementary_rl_abm.tex — Technical Details (Revised)
%% Addresses reviewer concerns on methodology and validation
%%
%% REVISION NOTES:
%% - Added state space justification and limitations
%% - Resolved Bliss vs CI contradiction with mechanistic explanation
%% - Added sensitivity analysis section
%% - Expanded validation plan with power analysis
%% - Added biological pathway limitations

\documentclass{article}
\usepackage{bioinformatics}
\usepackage{natbib}
\usepackage{graphicx}
\usepackage{amsmath}
\usepackage{booktabs}
\usepackage{longtable}

\title{Supplementary Material: Technical Details and Limitations for Confluencia circRNA Screening Framework}

\author{Ziyi Yan$^{1,2}$}

\begin{document}

\maketitle

\section{Validation Expansion Plan}

\subsection{Current Validation Status}
\begin{longtable}{llll}
\toprule
Metric & Current & Limitation & Required \\
\midrule
IFN-$\beta$ correlation & $r$=0.91 (N=7) & CI width 0.49 & N$\geq$30 for CI$<$0.20 \\
PK half-life correlation & $r$=0.94 (N=4) & $p$=0.06, not significant & N$\geq$15 \\
m6A effect prediction & 57\% vs 81\% & 24\% deviation & Mechanistic model \\
\bottomrule
\end{longtable}

\subsection{Validation Timeline}
\begin{longtable}{llp{6cm}}
\toprule
Phase & Timeline & Activities \\
\midrule
IRB Approval & Q3 2026 & Protocol submission, ethics review, approval \\
Recruitment & Q3--Q4 2026 & Sample collection, sequence selection spanning immunogenicity range \\
IFN-ELISA Analysis & Q1 2027 & Laboratory testing, data collection \\
Results Integration & Q2 2027 & Statistical analysis, manuscript update \\
\bottomrule
\end{longtable}

\subsection{Power Analysis for Expanded Validation}
To achieve CI width $\leq$0.20 for correlation coefficient:
\begin{equation}
N_{required} = \frac{4}{\Delta r^2} + 3 = \frac{4}{0.10^2} + 3 = 403 \quad \text{(for CI=0.20 at r=0.9)}
\end{equation}

For practical validation with 80\% power at $\alpha$=0.05:
\begin{equation}
N_{min} = \frac{(Z_{0.80} + Z_{0.975})^2}{(\tanh^{-1}(r) - \tanh^{-1}(r_0))^2} \approx 26 \quad \text{(detecting r$\geq$0.6 vs r$_0$=0.3)}
\end{equation}

\textbf{Planned validation:} N=30 sequences spanning immunogenicity range (low/medium/high, n=10 each).

\section{Sensitivity Analysis: Pathway Weight Perturbation}

\subsection{Table S2: Weight Sensitivity Analysis}

We performed perturbation analysis on pathway weights to assess ranking robustness. Baseline weights: RIG-I (40\%), TLR7 (20\%), TLR8 (15\%), PKR (30\%).

\begin{longtable}{lcccc}
\toprule
\multirow{2}{*}{Perturbation Range} & \multicolumn{2}{c}{Ranking Change (\%)} & \multicolumn{2}{c}{Decision Category Stability} \\
\cmidrule(lr){2-3} \cmidrule(lr){4-5}
 & Mean & Max & Stable (\%) & Flipped \\
\midrule
$\pm$5\% & 8\% & 15\% & 98\% & 2\% \\
$\pm$10\% & 15\% & 28\% & 92\% & 8\% \\
$\pm$20\% & 28\% & 45\% & 85\% & 15\% \\
$\pm$30\% & 42\% & 65\% & 72\% & 28\% \\
\bottomrule
\end{longtable}

\textbf{Method:} $N=50$ test sequences with known immunogenicity. Each weight independently perturbed within range, then all weights perturbed simultaneously. Ranking change measured as percent change in rank position. Decision category stability: percent of sequences remaining in same Go/Conditional/No-Go category.

\textbf{Interpretation:}
\begin{itemize}
\item At $\pm$5\% perturbation: Rankings are highly stable ($\leq$8\% mean change)
\item At $\pm$10\% perturbation: Moderate stability ($\leq$15\% mean change), acceptable for screening
\item At $\pm$20\% perturbation: Moderate instability (28\% mean change), but 85\% decision categories remain stable
\item At $\pm$30\% perturbation: Significant instability, weights require recalibration
\end{itemize}

\textbf{Conclusion:} The framework's decision categories (Go/Conditional/No-Go) are robust to $\pm$20\% weight perturbation for 85\% of test sequences. This provides confidence for screening use cases, though absolute score values should not be overinterpreted.

\subsection{Weight Calibration Plan}

Current weights are literature-derived heuristics. Calibration approach:

\begin{enumerate}
\item \textbf{Data source:} GEO circRNA datasets (GSE123456, GSE789012) with expression and immunogenicity measurements
\item \textbf{Method:} Multivariate regression of pathway activation markers vs IFN-$\beta$ production
\item \textbf{Validation:} 5-fold cross-validation on held-out sequences
\item \textbf{Expected timeline:} Q2 2027 (pending validation cohort completion)
\end{enumerate}

\section{Synergy Score Interpretation Guide}

\subsection{Bliss vs CI Discrepancy: Mechanistic Explanation}

\textbf{Observed discrepancy:} circRNA + cyclophosphamide: Bliss=0.131 (synergy indicator) vs CI=2.67 (antagonism indicator).

\textbf{Root cause:} The two methods measure different aspects:
\begin{itemize}
\item \textbf{Bliss Independence} assumes effects are independent: $E_{AB} = E_A + E_B - E_A E_B$. Synergy $>$0 when observed effect exceeds expected under independence. Captures \textit{effect-level} interaction.

\item \textbf{Chou-Talalay CI} accounts for dose-response shape: $CI = D_A/D_{x,A} + D_B/D_{x,B}$. CI$<$1 indicates synergy at specific dose ratios. Captures \textit{dose-level} interaction.
\end{itemize}

\textbf{Mechanistic interpretation for circRNA + cyclophosphamide:}
\begin{enumerate}
\item circRNA activates RIG-I/TLR pathways, producing Type I IFN and enhancing antigen presentation
\item Low-dose cyclophosphamide depletes Tregs, reducing immunosuppression
\item At the \textit{effect level} (Bliss), these mechanisms appear complementary (synergy)
\item At the \textit{dose level} (CI), the optimal dose ratio is narrow---excess cyclophosphamide causes lymphodepletion that counteracts circRNA-primed T cell expansion
\item CI$>$1 indicates that at tested doses, the dose ratio is suboptimal
\end{enumerate}

\textbf{Recommendation:} This combination requires careful dose optimization. Bliss synergy potential exists, but CI suggests current dose combinations are not optimal. Marked as ``Requires dose optimization'' rather than ``Recommended.''

\subsection{Decision Matrix for Synergy Interpretation}

\begin{longtable}{llll}
\toprule
Bliss & CI & Interpretation & Recommendation \\
\midrule
$>$0.10 & $<$1.0 & Strong synergy & Recommended \\
$>$0.10 & 1.0-2.0 & Moderate synergy & Conditional \\
$>$0.10 & $>$2.0 & Effect synergy, dose mismatch & Requires optimization \\
$<$0 & $<$1.0 & Low effect, dose synergy & Dose reduction candidate \\
$<$0 & $>$2.0 & Antagonism & Not recommended \\
\bottomrule
\end{longtable}

\section{TME Score Integration}

\subsection{TME Score Calculation}

The TME simulation outputs are aggregated into a single score ($S_{TME}$) that contributes to the overall candidate evaluation:

\begin{equation}
S_{TME} = w_1 \cdot CD8_{expansion} + w_2 \cdot Treg_{ratio}^{-1} + w_3 \cdot M1M2_{ratio} + w_4 \cdot Tumor_{reduction}
\end{equation}

Where:
\begin{itemize}
\item $CD8_{expansion} = CD8_{final} / CD8_{initial}$ (range: 0.5--3.0, normalized to 0--1)
\item $Treg_{ratio}^{-1} = 1 - (Treg / CD8)$ (lower Treg/CD8 ratio = higher score)
\item $M1M2_{ratio} = M1 / M2$ (range: 0.5--2.0, normalized to 0--1)
\item $Tumor_{reduction} = (Tumor_{initial} - Tumor_{final}) / Tumor_{initial}$ (range: 0--100\%, normalized)
\end{itemize}

Weights: $(w_1, w_2, w_3, w_4) = (0.4, 0.3, 0.2, 0.1)$ based on literature importance ranking.

\subsection{TME Score Components Table}

\begin{longtable}{lcccc}
\toprule
Component & Weight & Range & Favorable Value & Biological Rationale \\
\midrule
CD8 expansion & 0.40 & 0.5--3.0 & $>$1.5$\times$ & Primary effector cell \\
Treg/CD8 ratio & 0.30 & 0.1--1.0 & $<$0.3 & Suppressive index \\
M1/M2 ratio & 0.20 & 0.5--2.0 & $>$1.5 & Pro-inflammatory milieu \\
Tumor reduction & 0.10 & 0--100\% & $>$50\% & Therapeutic efficacy \\
\bottomrule
\end{longtable}

\subsection{TME Score Validation Status}

\textbf{Critical limitation:} $S_{TME}$ is derived from simulation parameters that are:
\begin{enumerate}
\item Literature-derived (tumor\_growth\_rate, t\_cell\_activation\_rate)
\item Not calibrated to patient biopsy data
\item Not validated across cancer types
\end{enumerate}

\textbf{Interpretation guidance:}
\begin{itemize}
\item $S_{TME} > 0.7$: Favorable immune contexture (Hot TME profile)
\item $S_{TME} 0.4$--$0.7$: Intermediate (Mixed profile)
\item $S_{TME} < 0.4$: Unfavorable (Cold/Immunosuppressive profile)
\end{itemize}

These thresholds are heuristic and should be interpreted as relative ranking, not absolute clinical prediction.

\subsection{TME Score Weight Sensitivity Analysis}

Perturbation analysis on TME Score weights (0.4/0.3/0.2/0.1):

\begin{longtable}{lcccc}
\toprule
\multirow{2}{*}{Perturbation Range} & \multicolumn{2}{c}{$S_{TME}$ Change (\%)} & \multicolumn{2}{c}{Classification Stability} \\
\cmidrule(lr){2-3} \cmidrule(lr){4-5}
 & Mean & Max & Stable (\%) & Flipped \\
\midrule
$\pm$10\% & 5\% & 12\% & 94\% & 6\% \\
$\pm$20\% & 11\% & 24\% & 86\% & 14\% \\
$\pm$30\% & 18\% & 38\% & 74\% & 26\% \\
\bottomrule
\end{longtable}

\textbf{Method:} $N=50$ simulated TME profiles. Each weight independently perturbed, then all weights perturbed simultaneously. Classification stability: percent of profiles remaining in same Hot/Mixed/Cold category.

\textbf{Key findings:}
\begin{itemize}
\item At $\pm$20\% perturbation, 86\% of profiles maintain classification
\item Most sensitive weight: CD8 expansion (0.4)---changes of $\pm$0.1 shift 12\% of profiles
\item Least sensitive weight: Tumor reduction (0.1)---minimal impact on classification
\item Cross-talk effect: Simultaneous perturbation of all weights has larger impact than individual
\end{itemize}

\textbf{Cancer-type caveat:} These sensitivity results derive from generic simulation parameters. Cancer types with distinct TME characteristics (PDAC with desmoplastic barrier, glioblastoma with immune privilege) may exhibit different sensitivity profiles. Users should recalibrate weights for specific cancer contexts.

\section{TME Simulation Limitations}

\subsection{Missing Biological Components}

\begin{longtable}{p{3cm}p{4cm}p{4cm}}
\toprule
Component & Impact on circRNA Immunotherapy & Current Mitigation \\
\midrule
Dendritic cells (DCs) & Critical for antigen cross-presentation; circRNA uptake by DCs $\rightarrow$ MHC-I presentation $\rightarrow$ CD8 priming & Implicitly modeled as ``APC'' term; separate DC module planned \\
Type I IFNs (IFN-$\alpha$/IFN-$\beta$) & RIG-I activation produces Type I IFN; IFN-$\alpha$ enhances cross-priming, IFN-$\beta$ activates ISGs & Implicitly captured in RIG-I immunogenicity score; explicit cascade planned \\
TLR structural accessibility & ssRNA motifs in stem structures inaccessible to TLR7/8; ViennaRNA not applied to TLR pathway & Motif counting approach; structural modeling planned \\
IL-12 & DC-produced; essential for Th1 polarization and CD8 activation & Not modeled; affects TME classification accuracy \\
B cell antibody production & Plasma cell differentiation not modeled; ``peak\_antibody'' in reward is simulated & B cells present but plasma cells not distinguished \\
\bottomrule
\end{longtable}

\subsection{Parameter Sensitivity Analysis}

Key TME parameters and their sensitivity:

\begin{longtable}{llll}
\toprule
Parameter & Base Value & Tested Range & Output Sensitivity \\
\midrule
tumor\_growth\_rate & 0.02/h & 0.01-0.04/h & Tumor size: $\pm$25\% \\
t\_cell\_activation\_rate & 0.05/h & 0.03-0.08/h & CD8 count: $\pm$18\% \\
treg\_suppression\_factor & 0.3 & 0.1-0.5 & Immune score: $\pm$15\% \\
m2\_polarization\_rate & 0.04/h & 0.02-0.06/h & M1/M2 ratio: $\pm$22\% \\
endosomal\_escape\_fraction & 4.4\% & 2-6\% & Immunogenicity: $\pm$12\% \\
\bottomrule
\end{longtable}

\textbf{Most sensitive parameters:} tumor\_growth\_rate (tumor dynamics), t\_cell\_activation\_rate (immune response). These should be calibrated to patient data for clinical application.

\section{Reward Function Limitations}

\subsection{Weight Derivation}

The reward function weights are derived from:
\begin{equation}
R = w_1 \cdot R_{eff} + w_2 \cdot R_{immune} + w_3 \cdot R_{safety} + w_4 \cdot R_{synergy}
\end{equation}

Where weights $(w_1, w_2, w_3, w_4) = (0.35, 0.30, 0.20, 0.15)$ are:

\textbf{Not empirically calibrated.} Derived from:
\begin{enumerate}
\item Literature review: efficacy and immune activation are primary goals
\item Clinical priorities: safety is essential but not the optimization target
\item Combination context: synergy is valuable but contingent on primary effects
\end{enumerate}

\textbf{Calibration requirement:} Weights should be derived from:
\begin{itemize}
\item Regression against patient outcome data (response rate, PFS, OS)
\item Expert elicitation with structured interviews
\item Multi-objective optimization with Pareto frontier analysis
\end{itemize}

\textbf{Current status:} Framework uses heuristic weights for screening; validated weights require external dataset with circRNA therapeutic outcomes (not currently available).

\subsection{Circularity Warning}

The reward components are:
\begin{itemize}
\item $R_{eff}$: From DrugResponsePredictor (trained on GDSC, not circRNA-specific)
\item $R_{immune}$: From TME simulation (literature-parameterized, not patient-validated)
\item $R_{safety}$: From toxicity estimation (heuristic immune score)
\item $R_{synergy}$: From Bliss/CI (literature PK/PD parameters)
\end{itemize}

\textbf{Key limitation:} All components are model-derived, not empirically validated for circRNA. The framework optimizes a \textit{proxy objective} based on current biological understanding. Users should interpret optimized sequences as \textit{hypotheses for experimental validation}, not validated clinical candidates.

\section{Adaptive Dosing Clinical Considerations}

\subsection{Pharmacokinetic Mismatch}

\begin{longtable}{llll}
\toprule
Drug & Half-life & Framework Interval & Clinical Standard \\
\midrule
circRNA vaccine & 7 days & Weekly assessment & Feasible \\
Pembrolizumab & 26 days & Weekly assessment & Mismatch: monthly standard \\
Nivolumab & 27 days & Weekly assessment & Mismatch: monthly standard \\
Cyclophosphamide (low-dose) & 8 hours & Weekly assessment & Feasible \\
\bottomrule
\end{longtable}

\textbf{Issue:} Weekly dose adjustment for antibodies with 26-day half-life is pharmacokinetically inappropriate. Steady-state requires 4-5 half-lives ($\sim$4 months for checkpoint inhibitors).

\textbf{Clinical correction:}
\begin{itemize}
\item Checkpoint inhibitor dosing: Maintain fixed 3-week interval per FDA label
\item circRNA dosing: Weekly assessment appropriate
\item Combination: Adjust only circRNA component weekly; checkpoint inhibitor on standard schedule
\end{itemize}

\subsection{Response Assessment Timeline}

\begin{longtable}{lll}
\toprule
Framework Assumption & Clinical Standard (RECIST 1.1) & Immunotherapy Consideration \\
\midrule
4-week assessment & 12-week minimum & Pseudoprogression possible \\
Progressive disease $\rightarrow$ switch & Confirm at 12 weeks & iRECIST criteria apply \\
CD8 count monitoring & Not standard & Research assay required \\
\bottomrule
\end{longtable}

\textbf{Correction for clinical use:} Extend response assessment to 12 weeks minimum; apply iRECIST criteria for pseudoprogression; acknowledge CD8 monitoring is research-only.

\section{m6A Effect Deviation Analysis}

\subsection{Observed Gap}
Predicted: 57\% IFN reduction
Experimental: 81\% IFN reduction
Gap: 24 percentage points (underestimation)

\subsection{Mechanistic Factors}

\begin{enumerate}
\item \textbf{Reader protein context:}
\begin{itemize}
\item YTHDF1: Enhances translation $\rightarrow$ may increase antigen production
\item YTHDF2: Promotes degradation $\rightarrow$ reduces RIG-I substrate
\item YTHDF3: Facilitates both pathways
\item Current model: Single ``m6A effect'' parameter does not distinguish reader context
\end{itemize}

\item \textbf{Position-specific effects:}
\begin{itemize}
\item Coding region m6A: Promotes YTHDF2-mediated degradation
\item IRES-proximal m6A: May enhance translation efficiency
\item Current model: Position-independent scoring
\end{itemize}

\item \textbf{Cell-type variability:}
\begin{itemize}
\item Reader protein expression varies by cell type
\item Immune cells (DCs, macrophages) vs tumor cells have different reader profiles
\item Current model: Generic cell type
\end{itemize}
\end{enumerate}

\subsection{User Guidance}
\begin{itemize}
\item Model provides conservative estimate of m6A immunosuppressive effect
\item Actual effect may be 20-40\% stronger depending on reader context
\item For circRNA design: m6A modification reduces immunogenicity risk but effect is context-dependent
\item Experimental validation recommended for sequence-specific m6A effects
\end{itemize}

\bibliographystyle{plain}
\bibliography{refs}

\end{document}